\begin{center}
    \large
    \begin{singlespace}    
        \textbf{\chineseTitle{}} \\[0.5cm]
    \end{singlespace}

    \begin{singlespace}    
    學生：\studentCnName  \hspace{2.5cm}  指導教授：\advisorCnName \hspace{0.1cm} 博士 \\

    \ifdefined\advisorCnNameB % 如果有共同指導教授
        \hspace{9.6cm}  \advisorCnNameB ~ 博士 \\
    \fi
    \end{singlespace}

    \vspace{0.5cm}

    \begin{singlespace}
        國立陽明交通大學電機學院\\
        \NameofDepartmentInstituteCN\\[0.2cm]
    \end{singlespace}

    \textbf{摘\hspace{1cm}要} \\[0.5cm]

\end{center}

\normalsize 

% 中文摘要就從這邊開始寫.
在過去文獻中，研究人員常使用基於梯度之伴隨方法(Adjoint Method)來設計光學結構。
但這種方法容易陷入局部最佳解。
因此，本研究應用因子分解機(Factorization Machines, FM)結合模擬退火(Simulated Annealing, SA)技術，
建構FMSA架構並應用於二維矽光子結構之反向設計(Inverse Design)，以探索全域最佳結構。

由於光學元件通常具有高維度設計空間，直接以二值像素進行最佳化將增加代理模型建構與搜尋的困難，
因此本研究進一步導入不同的編碼方式作為結構表示方法，來降低設計空間複雜度，同時提升結構表示能力與最佳化效率。

本研究具體採用了三種編碼策略，包括高斯上採樣(Gaussian Upsampling)、
二維傅立葉積分(2D Fourier Integral)以及變分自編碼器($\beta$-Variational Autoencoder,$\beta$-VAE)。

光學模擬的部分，本研究利用基於有限差分時域法(Finite-Difference Time-Domain, FDTD)的開源套件MEEP，
針對橫電(TE)模態在矽光元件中的傳輸行為進行模擬與分析，
並針對不同設計需求調整品質因子(Figure of Merit,FoM)以達不同設計目標。

本研究設計並分析了三種矽光子關鍵元件，
包括 彎曲波導(Bend Waveguide)、交叉波導(Waveguide Crossing)以及模態分波解多工器(Mode-Division Demultiplexer, DEMUX)。
其中，彎曲波導以提升穿透率(Transmission)為主要設計目標；交叉波導則利用結構對稱性降低設計變數維度，以有效抑制模態轉換與串擾(Crosstalk)，並維持高傳輸效率；
模態分波解多工器則以實現不同模態光場之分離為設計目標，提升各輸出通道之模態純度與傳輸效率。

透過 FM 模型結合模擬退火(FMSA)架構及適度調整FoM，成功設計出高效能矽光子元件，
並比較 Gaussian Upsampling、2D Fourier Integral 與 $\beta$-VAE 三種編碼策略於不同光學元件之適用性，
可作為未來矽光子反向設計與光子積體電路(Photonic Integrated Circuits, PICs)設計之參考。
\vspace{1cm}

% 中文摘要及關鍵詞 5-7 個 
\textbf{關鍵字：}因子分解機、矽光子、波導、反向設計、量子退火、機器學習
